\section{Introduction}
\label{intro}

Accurate short-term precipitation forecasting is important for urban flood management, transportation safety, and water-resource planning, particularly in regions such as Ireland where weather conditions can change rapidly under strong maritime influence \citep{toth2000comparison}. Precipitation remains difficult to predict because of the nonlinear evolution of atmospheric processes, the intermittency of rainfall, and limitations in the spatial and temporal representation of meteorological observations \citep{magnusson2019ecmwf}. Numerical Weather Prediction systems, including the European Centre for Medium-Range Weather Forecasts (ECMWF) Integrated Forecasting System, provide physically based forecasts and commonly use ensemble approaches to represent forecast uncertainty. \textcolor{red}{The present study evaluates a regional data-driven forecasting setting and does not claim that its models can replace operational NWP systems or provide a computational advantage over them.}

Machine-learning weather prediction has developed rapidly in recent years. \revminor{Global data-driven models such as FourCastNet, Pangu-Weather, GraphCast and FuXi illustrate different neural-network approaches to atmospheric prediction \citep{pathak2022fourcastnet,bi2023accurate,lam2023learning,chen2023fuxi}.} \textcolor{red}{Their reported performance and inference characteristics do not establish the accuracy or end-to-end computational cost of the regional models evaluated here.} Recurrent neural networks and related sequence models can represent temporal dependencies in sequences of meteorological predictors. However, predictive accuracy alone does not provide a complete assessment of precipitation forecasting performance. Rainfall distributions are highly skewed, dry and low-rainfall observations are common, and high-intensity events occur much less frequently. A model that performs well over the full test distribution may therefore behave differently when evaluated specifically under stronger rainfall conditions.

Uncertainty estimation provides an additional perspective on forecast quality. Many machine-learning forecasting models are trained to produce a single deterministic prediction, whereas risk-sensitive applications also require information about the uncertainty associated with that estimate \citep{yang2024interpretable}. Probabilistic neural networks can address this issue by predicting parameters of a conditional distribution rather than a single rainfall value. However, satisfactory probabilistic performance over the full test set does not necessarily imply reliable uncertainty estimates for less frequent heavy-rainfall events. Evaluation of prediction-interval coverage and probabilistic scores across different rainfall regimes is therefore important when assessing the reliability of probabilistic precipitation forecasts. In this study, these questions are evaluated for the upper tail of ERA5-derived regional-mean hourly precipitation, rather than as an assessment of operational heavy-rainfall warning capability.

Interpretability presents a related challenge. Deep neural networks can capture nonlinear temporal relationships but often provide limited information about how individual meteorological predictors contribute to their forecasts \citep{du2019techniques}. Attention mechanisms provide one way of weighting temporal representations, while post hoc attribution methods such as SHapley Additive exPlanations (SHAP) \citep{lundberg2017unified} can be used to examine predictor-level contributions. These interpretations should be distinguished from statistical association and predictive necessity. A variable may be strongly correlated with precipitation while providing limited additional information once closely related predictors are already included in the model. Combining correlation analysis, model attribution, and feature ablation can therefore provide a more complete assessment of predictor relevance.

Regional precipitation forecasting over Ireland provides a useful setting in which to examine these issues. Previous work has explored uncertainty-aware recurrent modelling for rainfall prediction in this region \citep{10906190}, while the rapidly changing Atlantic-influenced weather conditions continue to present challenges for localized forecasting. In this study, reference, machine-learning, deterministic deep-learning, and probabilistic recurrent models are evaluated for 6-hour-ahead prediction of hourly regional-mean precipitation over Dublin using ERA5 reanalysis data from 2004 to 2021. The target is the hourly-mean precipitation rate in mm\,h$^{-1}$ for the hour ending at the forecast valid time, not accumulated precipitation over the next 6~h. The analysis considers overall point-forecast accuracy, sensitivity to random initialization, performance across rainfall regimes, event-detection skill, probabilistic reliability, and predictor attribution. Particular attention is given to whether conclusions based on aggregate test performance remain valid for stronger precipitation events and whether different measures of predictor relevance lead to consistent interpretations.

\subsection{Main contributions}

The main contributions of this study are as follows:

\begin{itemize}
    \item A common experimental setting is used to compare reference forecasts, XGBoost, Transformer, deterministic BiLSTM models, and uncertainty-aware recurrent models for 6-hour-ahead prediction of hourly regional-mean precipitation over Dublin, rather than a 6-hour precipitation total.

    \item Model robustness is assessed using repeated training with five fixed random seeds, together with rainfall-regime and threshold-based analyses that distinguish overall performance from forecasting skill under the evaluated wet regime and upper-tail precipitation conditions.

    \item Probabilistic recurrent models are evaluated using negative log-likelihood \revminor{(NLL)}, continuous ranked probability score, prediction-interval coverage, and interval width. Their reliability is examined separately for the full test distribution and the upper tail of the precipitation distribution.

    \item Predictor relevance is examined using training-period correlation diagnostics and matched five-seed SHAP, single-variable ablation and group ablation in two architectures, distinguishing statistical association, regime-dependent model attribution and incremental or shared predictive utility.
    \item \textcolor{red}{A comparison of Gaussian, Tobit and Bernoulli--Gamma likelihoods within one probabilistic backbone assesses non-negative support, overall proper scores and remaining upper-tail undercoverage, while disclosing scale and dry-definition limitations.}
\end{itemize}

The remainder of this paper is organized as follows. Section~2 reviews related work on precipitation forecasting, machine-learning methods, probabilistic prediction, and model interpretability. Section~3 describes the study region, ERA5 dataset, feature selection, forecasting models, experimental setup, and diagnostic analyses. Section~4 presents the forecasting, robustness, probabilistic, and interpretability results. Section~5 discusses the main findings, their implications, and the limitations of the study. Section~6 summarizes the conclusions and directions for future work.

%Relatedwork----------------------------------
\section{Related work}

Recent advances in meteorological modeling have been supported by the increasing availability of high-resolution reanalysis datasets and the development of machine learning techniques, both of which have become important for precipitation forecasting, as demonstrated by \citet{10512739}. Reanalysis datasets provide temporally and spatially consistent climate information that supports retrospective climate studies and weather prediction applications \citep{10512483}. Among these products, the ERA5 reanalysis dataset developed by ECMWF under the Copernicus Climate Change Service represents a substantial improvement over its predecessors. \citet{hersbach2020era5} introduced ERA5 with hourly coverage and a spatial resolution of approximately 31~km, showing marked improvements over ERA-Interim in reproducing variability and extremes across multiple atmospheric variables. Subsequent evaluations have further clarified the reliability and limitations of ERA5. \citet{mcnicholl2021evaluating} found strong agreement between ERA5 near-surface air temperature and ground-based observations but also reported discrepancies under complex meteorological conditions. \citet{doddy2021reanalysis} compared several reanalysis products for renewable energy applications in Ireland and showed that ERA5 generally outperforms MERRA-2 and JRA-55, while still exhibiting regional biases that require consideration. For precipitation, \citet{lavers2022evaluation} evaluated 24-hour totals against 5,637 stations over 2001--2020, finding a general wet bias, lower annual maximum one-day amounts, and smaller random errors in winter than summer in the extratropics. These daily-scale findings do not establish the accuracy of hourly precipitation over Dublin. \revminor{The Copernicus European Regional Reanalysis (CERRA) provides atmospheric reanalysis data over Europe at a horizontal resolution of 5.5~km, offering finer spatial detail than the global ERA5 reanalysis \citep{ridal2024cerra}. Although \citet{ridal2024cerra} reported improvements over ERA5 for several near-surface variables, the present study uses ERA5 for both predictors and target precipitation and does not assess whether CERRA would improve regional precipitation forecasting.} \citet{palharini2021assessment} evaluated eleven satellite-based precipitation products across different regions of Brazil and found that no single product consistently captured extreme rainfall, with performance varying by region and event intensity. Similarly, \citet{kumar2025long} evaluated CHIRPS and GSMaP\_ISRO rainfall estimates against gauge observations over India during 2000--2022 and found that GSMaP\_ISRO agreed more closely with station measurements than CHIRPS, while both products diverged notably over complex terrain such as the Western Ghats. These studies show that precipitation products can provide useful regional information, but their performance remains dependent on location, terrain, and rainfall intensity.

\revminor{Parallel to advances in reanalysis data, machine learning has gained prominence in precipitation forecasting. Among conventional machine-learning approaches, \citet{rahman2022rainfall} developed a fusion-based system for smart-city rainfall prediction that combines multiple classifiers with fuzzy logic, reporting improvements in prediction accuracy in their study. With the development of deep learning, more sophisticated spatiotemporal architectures have been explored. \citet{adewoyin2021tru} proposed TRU-NET, which incorporates attention mechanisms and ERA5 inputs to produce high-resolution rainfall predictions over the UK. Their model outperformed U-Net baselines and produced competitive results against dynamical forecasts. In their UK daily-rainfall evaluation, the TRU-NET CC variant underpredicted observed amounts above 20~mm\,day$^{-1}$ more than IFS. This daily-scale comparison does not establish the behaviour of the models evaluated here for Dublin regional-mean hourly precipitation. At the global scale, \citet{singh2021deep} trained a residual U-Net on reanalysis data in a cubed-sphere framework and reported improved correlation skill over an operational numerical weather prediction system, while also emphasizing the need for hybrid integration and probabilistic evaluation. \citet{miller2021analysis} examined short-term and seasonal rainfall forecasts in Kenya and found substantial regional differences in forecasting skill, highlighting the importance of localized validation. More recently, \citet{latif2023assessing} compared statistical, machine-learning and deep-learning approaches to rainfall prediction and concluded that machine learning performed well across climates and timescales, with LSTM models among the most effective. These studies demonstrate the potential of machine learning and deep learning for precipitation forecasting, while also showing that aggregate performance may not fully represent model behaviour under different rainfall conditions.}

Reliable uncertainty estimation is important for the application of deep learning models in scientific forecasting. \citet{kendall2017uncertainties} introduced a probabilistic regression framework in which neural networks jointly predict the mean and variance of a Gaussian likelihood, allowing data-dependent uncertainty to be estimated directly from the model. \citet{9001195} proposed a framework that combines assumed density filtering and Monte Carlo dropout to model both data and model uncertainty without retraining, reporting improvements in calibration and robustness across perception and control tasks. \citet{klotz2022uncertainty} compared several probabilistic deep learning approaches for rainfall-runoff modeling and found that asymmetric Laplacian mixture networks produced more accurately calibrated predictions than Gaussian or dropout-based approaches, particularly under extreme flow conditions. \citet{mena2021survey} reviewed a range of uncertainty estimation methods, including variational inference, Monte Carlo dropout, ensembles, and Dirichlet prior networks, and discussed their trade-offs in calibration, interpretability, and computational cost. These studies show that probabilistic modeling can provide information beyond deterministic point forecasts, but they also indicate that forecast reliability depends on the assumed predictive distribution and on the part of the target distribution being evaluated. This is particularly relevant for precipitation, where high-intensity events are relatively rare and may exhibit different uncertainty characteristics from the full sample distribution.

Interpretability has also become an important consideration in data-driven forecasting. Although complex neural networks can achieve high predictive accuracy, their internal decision processes are often difficult to examine directly. \citet{chakraborty2017interpretability} provided an early systematic overview of interpretability research, distinguishing model transparency, functionality, and different forms of explanation. \citet{li2022interpretable} developed a unified taxonomy linking interpretation, interpretability, and trustworthiness, providing a general basis for evaluating explanations produced by machine learning systems. \citet{vinuesa2021interpretable} discussed the role of interpretable deep learning in sustainable development and emphasized the value of models that incorporate meaningful and understandable representations. \citet{temenos2023interpretable} developed an interpretable deep learning framework combining a compact CNN with SHAP-based explanations to provide both local and global feature attributions in remote-sensing classification. \citet{yoganand2026shap} proposed a SHAP-integrated LeNet model with transfer learning for flood prediction in Telangana, combining engineered feature selection with attention-based convolutional layers and using SHAP to explain model outputs. \citet{yang2024interpretable} reviewed recent developments in interpretable machine learning for weather and climate prediction, highlighting the use of feature attribution, attention mechanisms, and physics-informed approaches to examine relationships learned by data-driven models. Collectively, these studies show the value of interpretability methods for examining model behaviour and identifying predictors that contribute to forecasts. However, attribution should be distinguished from statistical association and from the predictive contribution measured by feature removal, particularly when meteorological predictors contain overlapping information.

These studies motivate a joint examination of overall predictive accuracy, precipitation-regime-dependent performance, probabilistic reliability, and predictor attribution within a consistent experimental setting. The present study evaluates six-hour-ahead forecasts of hourly ERA5-derived regional-mean precipitation over Dublin through model comparison, repeated-run robustness analysis, rainfall-regime and event-detection evaluation, probabilistic assessment, SHAP attribution, and feature ablation. The analysis examines how model performance and predictor relevance depend on the evaluated precipitation regime and diagnostic method; its conclusions are conditional on this regional reanalysis task and the tested configurations.

